\documentclass[exercice,a4]{thedocument}%
\usepackage[french]{babel}%
\usepackage{pstricks}%
\usepackage{enumitem}%
\startdatakeys
    \source{21GENMATNC1}
    \BO{2026}
    \cycle{4}
    \competencesDuSocle{CA,RA}
    \objectifsApprentissage{C4B2eb,C4B3ca}
    \motsCles{Pythagore,conversion,proportionnalité}
\enddatakeys
\begin{document}%

\medskip

BAI est un triangle rectangle en A tel que BA $= 210$~cm et AI $= 155$~cm.

\medskip

\begin{enumerate}
\item Déterminer la longueur BI au cm près.

\textbf{Rédiger la réponse en faisant apparaître les différentes étapes.}

L'immeuble de Joanne possède 15 vitres rectangulaires.

Chaque vitre a pour longueur 210 cm et pour largeur 155 cm.

Lors d'une préalerte cyclonique Joanne pose de l'adhésif sur les deux diagonales de chaque vitre de l'immeuble.
\begin{center}
\psset{unit=1cm,arrowsize=2pt 3}
\begin{pspicture}(6,6.8)
%\psgrid
\psframe(1.5,1.8)(4.5,6.3)
\psline(1.5,1.8)(4.5,6.3)\psline(4.5,1.8)(1.5,6.3)
\uput[d](3,6.3){$155$~cm}\uput[r](4.5,4.05){210 cm}
\rput(3,0.2){Une bande d'adhésif est assimilée à une diagonale du rectangle}
\rput(3,6.6){Schéma de la situation}
\psline{->}(3,0.6)(2.3,3)\psline{->}(3,0.6)(3.7,3)
\end{pspicture}
\end{center}

\item Justifier que Joanne a besoin d'environ $5,22$ m d'adhésif pour une vitre.
\end{enumerate}

Joanne a 7 rouleaux d'adhésif de $10$ m chacun.

\begin{enumerate}[resume]
\item A-t-elle assez d'adhésif pour toutes les vitres ? Justifier la réponse.
\end{enumerate}

\bigskip

\showthesource
\end{document}
